% Options for packages loaded elsewhere
\PassOptionsToPackage{unicode}{hyperref}
\PassOptionsToPackage{hyphens}{url}
\PassOptionsToPackage{dvipsnames,svgnames,x11names}{xcolor}
\documentclass[
  11pt,
  a4paper,
]{article}
\usepackage{xcolor}
\usepackage[margin=20mm]{geometry}
\usepackage{amsmath,amssymb}
\setcounter{secnumdepth}{-\maxdimen} % remove section numbering
\usepackage{iftex}
\ifPDFTeX
  \usepackage[T1]{fontenc}
  \usepackage[utf8]{inputenc}
  \usepackage{textcomp} % provide euro and other symbols
\else % if luatex or xetex
  \usepackage{unicode-math} % this also loads fontspec
  \defaultfontfeatures{Scale=MatchLowercase}
  \defaultfontfeatures[\rmfamily]{Ligatures=TeX,Scale=1}
\fi
\usepackage{lmodern}
\ifPDFTeX\else
  % xetex/luatex font selection
  \setmainfont[]{DejaVu Serif}
  \setsansfont[]{DejaVu Sans}
  \setmonofont[]{DejaVu Sans Mono}
\fi
% Use upquote if available, for straight quotes in verbatim environments
\IfFileExists{upquote.sty}{\usepackage{upquote}}{}
\IfFileExists{microtype.sty}{% use microtype if available
  \usepackage[]{microtype}
  \UseMicrotypeSet[protrusion]{basicmath} % disable protrusion for tt fonts
}{}
\makeatletter
\@ifundefined{KOMAClassName}{% if non-KOMA class
  \IfFileExists{parskip.sty}{%
    \usepackage{parskip}
  }{% else
    \setlength{\parindent}{0pt}
    \setlength{\parskip}{6pt plus 2pt minus 1pt}}
}{% if KOMA class
  \KOMAoptions{parskip=half}}
\makeatother
\usepackage{longtable,booktabs,array}
\newcounter{none} % for unnumbered tables
\usepackage{calc} % for calculating minipage widths
% Correct order of tables after \paragraph or \subparagraph
\usepackage{etoolbox}
\makeatletter
\patchcmd\longtable{\par}{\if@noskipsec\mbox{}\fi\par}{}{}
\makeatother
% Allow footnotes in longtable head/foot
\IfFileExists{footnotehyper.sty}{\usepackage{footnotehyper}}{\usepackage{footnote}}
\makesavenoteenv{longtable}
\setlength{\emergencystretch}{3em} % prevent overfull lines
\providecommand{\tightlist}{%
  \setlength{\itemsep}{0pt}\setlength{\parskip}{0pt}}
\usepackage{xeCJK}
\setCJKmainfont[AutoFakeBold=1.4]{Droid Sans Fallback}
\setCJKsansfont[AutoFakeBold=1.4]{Droid Sans Fallback}
\setCJKmonofont{Droid Sans Fallback}
\emergencystretch=3em
\setlength{\parindent}{0pt}
\xeCJKDeclareCharClass{Default}{"2013,"00B7}
\usepackage{newunicodechar}
\newunicodechar{ℱ}{\ensuremath{\mathcal{F}}}
\renewcommand{\contentsname}{目录}
\usepackage{bookmark}
\IfFileExists{xurl.sty}{\usepackage{xurl}}{} % add URL line breaks if available
\urlstyle{same}
\hypersetup{
  pdftitle={He--Kruczenski 2309.12402v3：理论推导与暂缓问题},
  colorlinks=true,
  linkcolor={Maroon},
  filecolor={Maroon},
  citecolor={Blue},
  urlcolor={Blue},
  pdfcreator={LaTeX via pandoc}}

\title{He--Kruczenski 2309.12402v3：理论推导与暂缓问题}
\usepackage{etoolbox}
\makeatletter
\providecommand{\subtitle}[1]{% add subtitle to \maketitle
  \apptocmd{\@title}{\par {\large #1 \par}}{}{}
}
\makeatother
\subtitle{既有推导与问题说明合并排版}
\author{}
\date{2026-09-11}

\begin{document}
\maketitle

{
\setcounter{tocdepth}{2}
\tableofcontents
}
\section{顺序主线：统一解析基及原参数联合问题}\label{ux987aux5e8fux4e3bux7ebfux7edfux4e00ux89e3ux6790ux57faux53caux539fux53c2ux6570ux8054ux5408ux95eeux9898}

\subsection{1.
先固定定义，不根据输出选择}\label{ux5148ux56faux5b9aux5b9aux4e49ux4e0dux6839ux636eux8f93ux51faux9009ux62e9}

本轮采用原文印出的两条手征残差不等式，各取四维 Euclidean
范数，epsilon=.002。保留打印 FESR 中心、四项 raw-absolute .002、原
uniform midpoint cutoff、FF squared caps 6e-5
和所有质量、耦合及凝聚输入。实际双密度 L4 的 B(M)=100{[}M²+M(M+1)/2{]}
作为已声明的数值正则化保留，不能称它是原2309明示参数。

新模型取完整 M 模式 sine-cardinal 密度及其精确 Cauchy
变换。它是原文未指定插值的一种明确实现，不与旧 PV
混合近似相混。没有直接施加 T0=0、FF endpoint
等式、弹性饱和或峰质量条件。所有散射约束仍为原生节点的
\textbar S\textbar≤1。

\subsection{2.
一个函数产生所有行}\label{ux4e00ux4e2aux51fdux6570ux4ea7ux751fux6240ux6709ux884c}

令 phi\_j=pi(j+1/2)/M，D\_nj=(2-delta\_nM) sin(n phi\_j)/M。正弦正交给 D
S=I，因此转换保留全部 M 方向。对 q\_j(z)=sum\_n D\_nj
z\^{}n，b\_j=tan(phi\_j/2)/M，有 q\_j(-1)=-b\_j。未减除 Cauchy 核为
q\_j(z)+b\_j。

设 R 为任意双密度矩阵，Q 为对称双密度矩阵，Q\_ii=Cii、Q\_ij=Cij/2。令 b
为上述列向量。把未减除振幅写成 z 模式的转换为

\[
\begin{aligned}
c&=T_0+b^T(\sigma_1+2\sigma_2)+b^T(2R+Q)b,\\
a&=D(\sigma_1+2Rb),\\
d&=D(\sigma_2+(R^T+Q)b),\\
U&=DRD^T,\qquad V=DQD^T.
\end{aligned}
\]

于是整个振幅为

\[
\begin{aligned}
A(s,t,u)={}&c+a^Tp(s)+d^T[p(t)+p(u)]\\
&+p(s)^TU[p(t)+p(u)]+p(t)^TVp(u),\\
p_n(s)={}&z(s)^n.
\end{aligned}
\]

同位旋由原交叉组合构造，部分波始终为 f=(1/4) integral P\_l
T。没有用不同的核分别替换物理和阈下取值。

\subsection{3.
已有解析数据的复用}\label{ux5df2ux6709ux89e3ux6790ux6570ux636eux7684ux590dux7528}

认证 registry 的原始行作用于 z\^{}n 模态系数；对称 tu 非对角基函数为
(z\_t\^{}n z\_u\textsuperscript{m+z\_t}m
z\_u\^{}n)/2。复用时实施上述转换的转置，得到完整 nodal C\_flat
行；不能把原行直接当 H，也不能丢掉 Nyquist 模式。

逐文件核对 registry 内 SHA256，按 canonical\_wave\_index
而非文件列表位置排布。所有旧代码仍保持非执行状态。新 provider
只读取数值区间数据。

独立检查采用旧 tip 全部3876系数，在预定 j=0,25,49 上比较 source
行积与当前 CardinalAmplitude 的严格独立角积分。九项全部相容，见
A1\_source\_audit/source\_audit.json。这确认坐标转换、1/4 投影、未乘
kappa 及 tu 对称因子；不转移旧 tip 的可行性。

\subsection{4.
自由列可以复用，但必须重新验证}\label{ux81eaux7531ux5217ux53efux4ee5ux590dux7528ux4f46ux5fc5ux987bux91cdux65b0ux9a8cux8bc1}

在原生节点，Im q\_j(s\_i)=delta\_ij，交叉道核为实。因此

\[
\operatorname{Im} f^0_0(s_i)=\tfrac32\sigma_{1,i}+\sigma_{2,i}+D^0_i(R,Q),
\qquad
\operatorname{Im} f^2_0(s_i)=\sigma_{2,i}+D^2_i(R,Q).
\]

其余分波的自由虚部列精确为零。代码先检查转换后的区间包含这些恒等值，再写成精确自由列；双谱列不这样处理。求解入口再验证这一布局，故
residual 消元的前提没有被移除。常数实部锚分别为5/2和1。

\subsection{5. 严格 IR
内点存在的构造}\label{ux4e25ux683c-ir-ux5185ux70b9ux5b58ux5728ux7684ux6784ux9020}

取
v\_j=sin(phi\_j)，T0=0，sigma1=sigma2=v/8，R=vv\textsuperscript{T，Q=vv}T/2。相应单核为
h(s)=1+z(s)，其色散密度 sin(phi)\textgreater0。物理交叉道 Legendre 投影

\[
J_l(s)=\tfrac14\int_{-1}^1 P_l(\mu)h(t)\,d\mu
\]

可写成正密度乘外区 Q\_l 的积分，故每个保留 l
都严格为正。对非零高波，虚部分别正比于 9 J\_l sin(phi)、3 J\_l
sin(phi)、J\_l sin(phi)（I=0、2、1）。两 S 波还具有正的单谱项。

因此有限所有盘可由足够小的共同正缩放满足： tau \textless{} min\_j 2 Im
f\_j/(kappa\_j \textbar f\_j\textbar²)。 同一缩放还能满足两 chi
球及实际双密度界。它证明存在一个新族的有限严格 IR
内点；浮点候选仍须重新检查，不能用径向缩放修复负虚部。

\subsection{6. 进入 UV 的次序}\label{ux8fdbux5165-uv-ux7684ux6b21ux5e8f}

先通过完整新 H 的 IR 内点检查，再 prepare-current 绑定同一 H。原生 FF 的
K+iI 不变，散射 Gram 使用新
H，所以电流的代数构造可复用，但必须重新保存模型身份。新联合见证通过之前，不做三点相移和分辨率扫描。

保存的 Float64 H
证书与实际解析行的区间核验分别报告。它们的差别是算子舍入误差，不再是两种不同的解析核。连续全能全自旋认证仍不作为原文有限原型的额外前置条件。

\subsection{7. 缺失分辨率的解析行：Gauss
余项而非经验阶数拟合}\label{ux7f3aux5931ux5206ux8fa8ux7387ux7684ux89e3ux6790ux884cgauss-ux4f59ux9879ux800cux975eux7ecfux9a8cux9636ux6570ux62dfux5408}

对未由认证数据覆盖的 M/L，直接投影同一有限多项式振幅。令
\(d=|s-4|\)，角变量换成 \(u=d(1-\mu)/2\)。物理区的交叉割线位于
\(u\le-4\) 和 \(u\ge d+4\)；阈下 \(0<s<4\) 时位于 \(u\le-s\) 和
\(u\ge d+s\)。因此分别取 gap=4 或 gap=s，不能混用。

把 {[}0,d{]} 分成从端点向中部几何增长的对称面板 {[}l,r{]}，中心 c、半宽
h。令 D=min(gap+l,gap+d−r)，在面板标准坐标中选 Bernstein 椭圆的实半轴 a
满足 h(a−1)≤D/2。代码取向下舍入的 a≤1+D/(2h)，再取向下舍入的
\(\rho=a+\sqrt{a^2-1}>1\)。整个椭圆留在两个割线之间，故交叉 conformal
map 均满足
\textbar z\textbar≤1。于是任意保留幂次及其两个交叉道乘积的模均≤1，这个界不依赖
M。

若面板上 \(|\mu|\le R\)，Legendre 积分表示给出 \[
|P_\ell(\mu)|\le\bigl(R+\sqrt{R^2+1}\bigr)^\ell=:B_\ell.
\] 对椭圆内解析且模≤B的函数，Chebyshev 系数满足
\(|a_k|\le2B\rho^{-k}\)。用次数 2N−1 的截断多项式，余项一致界为
\(2B\rho^{-2N}/(1-\rho^{-1})\)。N点 Gauss
公式对该多项式精确，积分测度和正 Gauss 权重的总质量均为2，故 \[
\left|\int_{-1}^1g-Q_Ng\right|
\le {8B\rho^{-2N}\over1-\rho^{-1}}.
\] 本论文分波还含 1/4，角换元给 dμ=2du/d，面板换元给 du=h
dx；所以代码每面板加的是 \[
{4h\over d}\,{B_\ell\rho^{-2N}\over1-\rho^{-1}},
\] 而不是直接照搬8B。这给每个单模式积分 J 和双模式积分 U
的显式误差球。之后的同位旋组合、DST转置和未减除变换全部用区间矩阵运算传播误差，最后
Float64 H 的 radius\_upper
还包含存储舍入误差。N≥64是在看新曲线之前固定的计算选择，不是对曲线调参。

对称积分恒等式 Uji=(−1)\^{}ell Uij
在同一换元下成立；奇数ell时对角为零。只有原包络包含0时才写成精确零。原生自由虚部列也只在包络包含已证明的恒等值后精确化。独立角积分测试覆盖
S、P 以及非主分波的同一完整随机系数振幅。

一般精确/区间能量不按 float 相等关系吸附到原生节点；只有 CLI
明确约定的浮点节点别名或显式原生节点对象才使用缓存。非有限能量和跨越阈值的区间拒绝。这些接口条件保证``原生节点缓存''和``一般能量包络''不会互相冒充。

\subsection{8.
中心收敛、支持误差与物理曲线是三个问题}\label{ux4e2dux5fc3ux6536ux655bux652fux6301ux8befux5deeux4e0eux7269ux7406ux66f2ux7ebfux662fux4e09ux4e2aux95eeux9898}

对固定可行域，写障碍目标 \[
\begin{aligned}
\Psi_\mu(z)&=\Phi(z)-c^Tz/\mu,\\
\Phi&=-\sum_j\log\Delta_j-w_\chi\sum_a\log(\epsilon^2-\|r_a\|_2^2)
+\Phi_\rho(\rho/B)+\Phi_{\rm current}.
\end{aligned}
\] 这里实现的是保留原L4可行域的辅助变量消元障碍，不能误写成简单的
−log(1−sum r\_i\^{}4)。令 r\_i 为实际密度除以 B， \[
\Phi_\rho(r)=\min_{v_i>r_i^2,\;v^Tv<1}
\{-\sum_i\log(v_i-r_i^2)-\log(1-v^Tv)\}.
\] 存在这样的 v 当且仅当 sum r\_i\^{}4\textless1。记
a\_i=r\_i²、d\_i=v\_i−a\_i、δ=1−sum v\_i²，内层驻点满足 δ=2v\_i
d\_i，因此 \[
v_i={a_i+\sqrt{a_i^2+2\delta}\over2},\qquad
F(\delta)=\delta+\sum_i v_i(\delta)^2-1=0.
\] F(0)=sum a\_i²−1\textless0，F(1)\textgreater0，且 F′=1+sum
v\_i/sqrt(a\_i²+2δ)\textgreater0，故唯一正根存在。代码的标量Newton步正是对这个F求根，不是另设密度正则器。消元后的梯度为
g\_i=2r\_i/d\_i=4r\_i v\_i/δ。再隐式微分可得 \[
\nabla^2\Phi_\rho=\operatorname{diag}
\left({2\over d_i}+{8a_i\over\delta+2d_i^2}\right)+uu^T,
\quad
u_i={-4r_i v_i\over(\delta+2d_i^2)\sqrt{1+\sum_j\delta/(\delta+2d_j^2)}}.
\]
末项确实是正的秩一项。对完整辅助Hessian作Schur消元不能据直觉把最后的显式表示改成负号。恢复未归一化密度坐标时再乘1/B或1/B²；代码的特征矩阵与独立有限差分检查使用相同规则。

每个权重均为正。把 wχ 从1变为750不改变任何可行不等式，也不改变投影目标
c；它改变有限μ时选择的中心路径。选取750的规则是
n\_disks/2，与已使用的联合求解器一致，不由相移决定。梯度、Hessian、沿射线导数和对偶权重必须同时乘相应权重，不能只改停止量。

固定 \(Xz=x_{\rm ref}\) 时，以 \(V=(-X_{1:}/X_0;I)\) 参数化切空间（这里
X0≠0）。Newton 方程为 \[
(V^T\nabla^2\Psi V)d=-V^T\nabla\Psi,
\qquad \Delta z=Vd,
\] 因此全部密度方向仍在该仿射截面中优化。法向乘子由完整梯度沿 X
的分量恢复，形成全局支持方向
\((\lambda,n_y)\)。径向中心换元不能在这条固定截面迭代中直接缩放 z。

这里的联合中心是消去自由高能电流谱后的约化问题中心，仍保留全部3876个振幅方向。原始4076变量中的14个高能谱对角元既无矩上界也不进目标；若对它们直接使用完整−logdet障碍，沿R→∞会下降，不能宣称原始全变量问题有唯一有限障碍中心。严格散射内点上用Schur公式消去这些自由谱，再将约化中心拓展为满足全部原式的完整点，才是当前算法与证书的准确关系。

固定μ\textgreater0时，这个约化中心在严格内点域中是唯一的。上面的密度Hessian具有严格正对角加半正定秩一项，因此任何非零密度方向都有正曲率。若剩余自由方向使散射障碍的Hessian二次型为零，则所有原生分波的实、虚部变化都为零；先由S2的自由虚部列得δσ2=0，再由S0得δσ1=0，最后由S0实常数列5/2得δT0=0。因此IR
Hessian正定，固定xref的切空间限制也仍正定。

对联合约化问题，先由上述论证得δC=0。低能正定Gram的−logdet曲率为 tr(G⁻¹δG
G⁻¹δG)，为零当且仅当δG=0，因而所有低能ImF及受矩约束的R变化为零。高能FF盘的正二次曲率又要求各高能ImF变化为零。故联合约化Hessian也正定。密度/散射界使C有界，正FESR权重和FF上界使保留电流变量有界；严格内点已由新Phase
I实际获得，障碍在相应边界发散，所以每个固定μ的最小值存在且唯一。这解释了数值热启动不构成按相移选取另一个中心；它不证明μ→0的极值幅度唯一，也不证明ρ极点或物理曲线误差。

收敛中心要求 Newton
decrement、实际线性残差和残余梯度的支持宽度共同达标。早先``一次Newton后目标gap已小''只说明那一候选的投影间隙，不能证明它已到中心。本轮
\texttt{-\/-require-center}
保留真正收敛的完整中心作为代表，停止不再只由旧对偶和新非中心primal之间的gap触发。

已完成 tip 的6次精修给
decrement≈1.82e−14、支持gap≈7.44e−5。ref/mid也分别在其固定截面上完成中心与支持。最后两层μ间，ref/mid
的 P1 最大模π相位变化约0.256°/0.611°；这只是实测变化，不能当成 μ→0
的误差定理。支持gap只控制线性目标，不控制所有分波或保证极值振幅唯一。末端极小的浮点可行恢复缩放、以及自由高能电流谱的Schur提升均另行记录，不能隐称保存点与实数中心逐位相等。

IR
的权重1路径曾在极薄手征余量处停滞；直接重算状态发现最大分波坐标漂移约2.3e−14，手征余量漂移约3.6e−20，不支持``缓存状态明显漂移''的解释。采用相同750权重规则后，绿色IR中心收敛，支持gap约5.44e−4，对应纵向gap约1.09e−6；它又通过全部1500解析散射行、两条χ和L4的区间核验。

\subsection{9. 本轮已经闭合的 IR→UV→ρ
证据链}\label{ux672cux8f6eux5df2ux7ecfux95edux5408ux7684-iruvux3c1-ux8bc1ux636eux94fe}

在固定 xref
的绿色IR上侧中心，P1相移仍很小，未出现90°上穿或具有ρ强度的共振峰。完整C先保存于
C2\_green\_selection，之后才求值。三个UV代表同样在选点之前固定目标和规则，全部4076原变量的原约束及解析primal核验均通过。其90°读数为
tip约803.49、mid约701.26、ref约696.63MeV；原生散射强度主峰分别为792.14、680.41、680.41MeV。没有加入ρ质量、宽度或相位拟合条件。

固定绿色IR的整个C，只求电流扩展，原保存有限算子的Farkas上界约−9.10961718\textless0。证明形式是：所有电流Gram
PSD、FF锥和四个矩区间的非负组合，加上已由这些原约束推导的变量残差上界，推出零目标≤负数。PSD对偶修正经Arb主子式检查，不能只依赖求解器的
PrimalInfeasible
字样。此证书排除的是这个完整C，而不是其二维投影、整个联合可行域或所有``没有ρ''的函数。

Gram Schur 恒等式（\(\Delta=1-|S|^2>0\)）为 \[
R\ge |\mathcal F|^2+{|\mathcal F-S\mathcal F^*|^2\over\Delta}.
\]
F(0)=1与高能FF界限制形状因子的解析变化，UV矩限制可用电流谱预算，偏离散射/FF相位一致性的项消耗额外谱预算。因此原UV信息可以排除IR代表并推动出现共同中能结构。上述严格不等式没有自动要求Gram秩1，也没有证明唯一ρ极点或正确质量；矩约束本身也不是整个谱的唯一性条件。

所以本轮已复现``加入UV信息后出现ρ式结构''的有限数值机制；三代表中两条共振偏低且非弹性较强的
major
差异仍在。中心精修和共同解析原型修复没有消除这一差异，不能将其重新归类为小的Newton未收敛误差。新M/L、完整区域及收缩非对称仍须按顺序检查，不能用旧PV状态代替。

固定C证书现已用768-bit重新计入全部相关解析源行的误差球，并改用精确原生K、运动学、打印十进制矩、raw容差和hard权重。新严格上界为
−9.109617136909641，仍明确为负；相对原有限H上界仅变化约4.12e−8。证据见
CD\_fixed\_current\_analytic/source\_audit.json。这里的排除已不依赖忽略H的存储/积分误差。

\subsection{10.
新增的节点间反例}\label{ux65b0ux589eux7684ux8282ux70b9ux95f4ux53cdux4f8b}

本轮后续直接求值和18项独立积分发现三UV振幅在ρ附近有严格η\textgreater1，详见
\href{HELD_INTER_NODE_ISSUE_ZH.md}{独立汇总}。第9节的机制结论严格限于原生有限程序；它不能被解释为已经得到连续物理上合格的ρ相移。按照用户对这类反例的要求，网格加密暂列独立修补方案，未并入原基线。

\subsection{11.
两条向量矩如何容许偏低的ρ}\label{ux4e24ux6761ux5411ux91cfux77e9ux5982ux4f55ux5bb9ux8bb8ux504fux4f4eux7684ux3c1}

原文Eq.(3.72--74)写原始离散积分与QCD值的绝对差；Eq.(2.56)打印的是除以s0\^{}(n+2)后的数值，必须先乘回。当前逐项绝对容差.002有印式依据；原作者是否另做范数打包或代码缩放仍未公开，不能因为别的口径更接近相位就改用它。原文见
\texttt{references/2309.12402v3-source/prd\_submission\_2.tex:1022–1040}。

令正的离散测度权重为 a\_i=(π/M)s′\_i，保留s\_i≤s0。向量谱R\_i≥0，因此 \[
\begin{aligned}
I_{-1}&=\sum_i a_iR_i/s_i>0,\qquad
\mu_i={a_iR_i/s_i\over I_{-1}},\quad\sum_i\mu_i=1,\\
\overline{E^2}&=m_\pi^2\sum_i\mu_i s_i=m_\pi^2{I_0\over I_{-1}}.
\end{aligned}
\] 这里是以R/s为权重的能量平方均值，不是任意峰位或极点质量的公式。

打印中心值转回原始矩后是
I−1=0.004228045714\ldots、I0=0.145711206997\ldots。同样的绝对.002误差，对前者相当于47.3\%，对后者约1.37\%。只由这两条区间推出的加权均方能标范围约为0.6725--1.1399GeV，中心值为0.8219GeV。这既给出为何UV引入一个中能尺度，也说明它单独并不固定770MeV，更不保证谱有窄峰。

当前三个UV点实际使用的向量矩如下：

{\def\LTcaptype{none} % do not increment counter
\begin{longtable}[]{@{}lrrrr@{}}
\toprule\noalign{}
点 & I−1 & I0 & 加权均方能标 MeV & 直接加点90°读数 MeV \\
\midrule\noalign{}
\endhead
\bottomrule\noalign{}
\endlastfoot
tip & .00443308288 & .14771073369 & 808.13 & 802.01 \\
mid & .00590775693 & .14771083475 & 700.04 & 697.13 \\
ref & .00609361320 & .14771082229 & 689.28 & 688.38 \\
\end{longtable}
}

三点的I0均靠近允许上界，I−1的变化显著改变了这个谱能标。数据来自
\texttt{CE\_phases/phases.json}；没有用上述比值反向设定任何振幅或峰。它解释了为何当前低峰仍通过给定的UV矩，而不是把它归为Newton误差。

Gram条件进一步限制相位：R≥\textbar ℱ\textbar²+\textbar ℱ−Sℱ*\textbar²/(1−\textbar S\textbar²)。若IR的S几乎为1，而由归一化、解析性和高能界产生的F有较大虚部，则相位失配会消耗很大的有限谱预算；固定IR振幅的负Farkas界是这一冲突的全变量证据。改变S的相位或增加非弹性都可降低这项代价，原不等式没有强制二π饱和。

在δ=90°且F相位一致时，R的Schur下界为2\textbar ℱ\textbar²/(1+η)，对应二π谱比例最多达到(1+η)/2；散射强度则为(1+η)²/4。因此一个很高的FF/电流谱峰可以伴随较弱、非弹性很强的散射峰。当前原生三点的η差异正需要保留为major物理差异，不能用``Gram近饱和''代替弹性或完整ρ复现。

\clearpage

\section{节点间幺正性：已证实的缺陷与暂缓修改方案}\label{ux8282ux70b9ux95f4ux5e7aux6b63ux6027ux5df2ux8bc1ux5b9eux7684ux7f3aux9677ux4e0eux6682ux7f13ux4feeux6539ux65b9ux6848}

结论：本轮三个UV代表是原生有限程序的合法解，但不是在所检查区间处处满足幺正性的物理振幅。中心已经收敛，这个差异不能再归因于一次未完成的Newton更新。以下是当前已定义解析族的严格数值反例，不是论文有限SDP自身矛盾、作者未公开实现失败或原物理参数下整个联合问题不可行的证明。

\subsection{1.
冻结输入与结果}\label{ux51bbux7ed3ux8f93ux5165ux4e0eux7ed3ux679c}

三个完整振幅先在 \texttt{E4\_selection/\{tip,mid,ref\}}
保存，再直接求值。其原生1500个散射约束、100个电流Gram、两条χ、L4、四矩和高能FF界已经分别在
\texttt{E1\_tip\_verified}、\texttt{E2\_ref\_verified}、\texttt{E3\_mid\_verified}
通过解析误差核验。

新增检查保留这些系数，在0.60--0.90GeV加入每10MeV一个能量点；没有重新优化，没有剪裁η，没有以实验或论文相位选点。完整数据见
\texttt{E5\_direct\_profiles/direct\_profiles.json}，图见同目录
\texttt{direct\_profiles.pdf}。

{\def\LTcaptype{none} % do not increment counter
\begin{longtable}[]{@{}
  >{\raggedright\arraybackslash}p{(\linewidth - 8\tabcolsep) * \real{0.1579}}
  >{\raggedleft\arraybackslash}p{(\linewidth - 8\tabcolsep) * \real{0.2105}}
  >{\raggedleft\arraybackslash}p{(\linewidth - 8\tabcolsep) * \real{0.2105}}
  >{\raggedleft\arraybackslash}p{(\linewidth - 8\tabcolsep) * \real{0.2105}}
  >{\raggedleft\arraybackslash}p{(\linewidth - 8\tabcolsep) * \real{0.2105}}@{}}
\toprule\noalign{}
\begin{minipage}[b]{\linewidth}\raggedright
代表
\end{minipage} & \begin{minipage}[b]{\linewidth}\raggedleft
原生90°读数 MeV
\end{minipage} & \begin{minipage}[b]{\linewidth}\raggedleft
加点后描述读数 MeV
\end{minipage} & \begin{minipage}[b]{\linewidth}\raggedleft
加点最大η
\end{minipage} & \begin{minipage}[b]{\linewidth}\raggedleft
严格违反的波/能量样本数
\end{minipage} \\
\midrule\noalign{}
\endhead
\bottomrule\noalign{}
\endlastfoot
tip & 803.4866 & 802.0117 & 1.0925161 & 42 \\
mid & 701.2588 & 697.1341 & 1.1189324 & 40 \\
ref & 696.6313 & 688.3834 & 1.1409410 & 40 \\
\end{longtable}
}

两列90°读数都是插值描述，不是极点质量或严格根包络。新增求值并没有把ref/mid共振移回正确位置。

在0.70和0.90GeV对应的实际输入浮点能量上，又使用完整振幅的独立
rational-circle Arb
积分复算全部三主波、三代表，共18项，全部与Gauss解析行求值相容；其中11项严格违反幺正性。证据在
\texttt{E5\_independent\_integrals/source\_audit.json}。例如tip的P1在0.90GeV附近满足

\[
1-|S_1^1|^2\in[-0.1935914174659331795334\pm9.17\times10^{-23}]<0.
\]

因此它不是只由网格连线产生的视觉波纹，也不是未计入积分误差所致。首次窗口求值因自动电流诊断仅支持原生节点而未写出evaluation；其已完成的违例数据和失败报告仍保留。接口已改为明确标注该电流诊断不适用于离节点网格，散射求值和违例数据正常保存，未修改任何物理系数。

\subsection{2.
反例准确否定什么}\label{ux53cdux4f8bux51c6ux786eux5426ux5b9aux4ec0ux4e48}

设 K\_native 是本轮固定M50/L10、固定全部原输入的原生有限可行集；K\_all
是同一解析族同时在所有物理能量和保留波上满足幺正性的集合。显然
K\_all⊆K\_native。上述完整点证明这个包含可以严格：它属于K\_native，却不属于K\_all。

解析性不能补上这个缺口。最大模原理需要整个适当边界上的界，有限个边界点上的模界不够。密度L4界使固定M的系数有界，也没有把有限点检查变成区间上界；其数值大小更不提供这种自动结论。提高Arb位数只会更精确地识别这些违例。继续求同一个有限问题的更小μ中心也没有满足未加入约束的保证。

这不撤销已证明的有限IR→UV机制：无ρ的绿色IR完整振幅被UV电流扩展排除，三个UV代表出现ρ式结构，证书与原生峰都仍成立。但``得到正确的ρ及物理上合格的相移''属于更强判断，目前不成立。这应保留为major问题，不能把η剪裁为1或把所画曲线称为已验证的连续物理解。

\subsection{3.
为什么先不改原问题}\label{ux4e3aux4ec0ux4e48ux5148ux4e0dux6539ux539fux95eeux9898}

用户要求：对学术上自洽的反例或论文没有解决的问题，另行汇总并先不要修改。原文有限网格没有提供连续认证；作者散射插值与数值正则化也没有唯一公开。因此这里保留原文参数下的有限基线，并继续其区域及五组M/L比较，不把新增约束悄悄并入它。这个暂缓只针对改变数值约束集合；求值、诊断、模型身份或中心停止逻辑中的实现错误仍正常修复。

\subsection{4.
可行的下一版修补}\label{ux53efux884cux7684ux4e0bux4e00ux7248ux4feeux8865}

最直接的方案是保持M、L、全部密度方向、物理参数、χ范数、FESR和FF输入不变，将散射检验网格与基函数模式数分开：

\begin{enumerate}
\def\labelenumi{\arabic{enumi}.}
\tightlist
\item
  保留原生节点及其自由列恒等式，在同一解析族上生成额外物理能量的分波行。
\item
  先加入已由独立计算严格确认违反的能量点，重新求同一投影目标，始终保持\textbar S\textbar≤1。
\item
  在预先固定的更细能量网格复查；有实际违例再添加约束。这个交换过程由物理残差驱动，不由ρ目标质量或实验曲线驱动。
\item
  所有IR和UV代表重新选定、重新验回并贯穿对应图；不把新点的可行性与旧点的好看相位拼接。
\item
  输出新旧有限程序的支持区间、峰、相移、η和采样收敛。附加有限网格通过仍不称完整连续认证；若确有必要，可在相关有限能区补区间上界，而不要求先证明无限能量/无限自旋问题。
\end{enumerate}

实现上已有一般能量解析行、严格IR种子、完整联合Phase
I和原式对偶核验。需把prepare/audit/cache的行布局改为显式逐行能量/分波元数据，保留原生行在前，更新相应签名。它不需要改变底层物理约束，也不需要导入论文答案；但会产生一套更密配点的新有限程序，因此本轮先列为独立提案。

\subsection{5.
原五组分辨率能回答什么}\label{ux539fux4e94ux7ec4ux5206ux8fa8ux7387ux80fdux56deux7b54ux4ec0ux4e48}

同M时L8→L10→L12只增加保留分波；若共同数值行逐项一致，其有限可行集应嵌套，+x支持上界应相应收缩。不同M同时改变函数空间、能量节点、FESR求积和声明的B(M)，并非一串嵌套问题，不能用``变量更多''推断单调性。原Fig.11比较是稳定性证据，不是连续极限定理。五组都应继续使用预先规定的+x代表；不得通过重新选择相移较好的点隐藏本次节点间违例。

\subsection{6.
同一有限解析族还有一个必要的无穷远条件}\label{ux540cux4e00ux6709ux9650ux89e3ux6790ux65cfux8fd8ux6709ux4e00ux4e2aux5fc5ux8981ux7684ux65e0ux7a77ux8fdcux6761ux4ef6}

在当前未减除节点密度坐标中，h\_j(s)=q\_j(z(s))+b\_j，且q\_j(−1)=−b\_j，所以h\_j(s)在s→∞时趋于0。对物理散射角的几乎所有μ，s→+∞时t,u→−∞，整个有限多项式振幅趋于第一系数T0。各conformal变量在闭单位盘上，有限多项式有统一有限界，因此可用控制收敛定理交换角积分和极限。得到
\[
f_0^0(s)\longrightarrow\tfrac52T_0,
\qquad f_0^2(s)\longrightarrow T_0,
\qquad \kappa\longrightarrow\pi.
\] T0为实数，所以全能量幺正性的必要条件是 \[
1+(5\pi T_0/2)^2\le1,\qquad 1+(\pi T_0)^2\le1,
\quad\Longrightarrow\quad T_0=0.
\]
这针对当前有限多项式密度的无穷远行为；不把它外推给具有不同密度尾部的所有振幅，也不声称原作者使用了同一参数化。

当前tip/mid/ref的T0分别为27.61688919、−70.71003327、−64.93454399，对应S0的η极限约216.90、555.36、510.00。\texttt{E5\_direct\_profiles\_complete/direct\_profiles.json}保存了这些极限的Arb包络。因而这些完整C不可能在所有能量上幺正，独立于ρ窗口中已经发现的违例。

这里的零条件在减除坐标中对应一个线性组合等式；低能接触强度仍可由谱积分产生。它也不是充分条件：设T0=0之后，节点间仍须检查。对当前有限形式因子F=1+sum
q\_j
ImF\_j，同样有F(∞)=1−b·ImF；向量k1²随s线性增长，因此把原FF上界解释为所有s\textgreater s0上的界时，还必须有F1(∞)=0。数值采样通过并未自动保证这些端点条件。

T0=0这一必要切片在当前原生联合问题中至少有可行点：tip的T0为正、ref为负，取
\[
t={-T_{0,\mathrm{ref}}\over T_{0,\mathrm{tip}}-T_{0,\mathrm{ref}}}\approx0.70160495
\]
混合两个完整联合点，则T0精确为0。散射盘、χ球、密度L4球、FF盘、FESR区间及Gram
PSD全都凸，F(0)=1是仿射归一化，所以这个精确凸组合保持全部原生原约束。其x约0.09137。这只证明该有限切片非空，没有证明混合点的节点间幺正性或正确ρ；本轮没有以它替换任何代表或相移。

据此，最小的后续修补提案可先在同一有限族内加入已推导的必要端点关系，再加密散射检查；若希望允许不同的高能吸收行为，则需先明确并推导新的密度尾部基函数。两条路线的函数族及近似含义应分别声明。本轮依用户要求把这个解析反例另行汇总，保持原生基线不变。

\end{document}
